\documentclass{article}
\usepackage{graphicx}
\usepackage{hyperref}
\usepackage{fancyhdr}
\usepackage{indentfirst}
\usepackage{graphicx}
\usepackage{newlfont}
\usepackage{amssymb}
\usepackage{amsmath}
\usepackage{latexsym}
\usepackage{lipsum}
\usepackage{tikz}
\usepackage{amsthm}
\usepackage{algorithm}
\usepackage{algpseudocode}
\usepackage{stmaryrd}
\usepackage{listings}
\usepackage{xcolor}
\usepackage{xcolor, colortbl}
\usepackage{pgfgantt}
\usepackage{graphicx}
\usepackage{dirtree}
\usepackage{todonotes}
\usepackage{bussproofs}
\usepackage{mathtools}
\usepackage{msc}
\usepackage{subcaption}

\newtheorem{theorem}{Theorem}
\newtheorem{lemma}{Lemma}
\newtheorem{corollary}{Corollary}

\usetikzlibrary{automata,arrows,positioning}
\pgfmathtruncatemacro\distance{1}

\definecolor{green}{rgb}{0,0.5,0}
\definecolor{red}{rgb}{1,0,0}
\definecolor{yellow}{rgb}{0.5,0.5,0}
\definecolor{codeashgrey}{rgb}{0,0.6,0}
\definecolor{codegray}{rgb}{0.5,0.5,0.5}
\definecolor{codepurple}{rgb}{0.58,0,0.82}
\definecolor{backcolour}{rgb}{0.95,0.95,0.92}
\definecolor{ashgrey}{rgb}{0.7, 0.75, 0.71}

\lstdefinestyle{mystyle}{
    backgroundcolor=\color{backcolour},   
    commentstyle=\color{codeashgrey},
    keywordstyle=\color{magenta},
    numberstyle=\tiny\color{codegray},
    stringstyle=\color{codepurple},
    basicstyle=\ttfamily\footnotesize,
    breakatwhitespace=false,         
    breaklines=true,                 
    captionpos=b,                    
    keepspaces=true,                 
    numbers=left,                    
    numbersep=5pt,                  
    showspaces=false,                
    showstringspaces=false,
    showtabs=false,                  
    tabsize=2
}

\lstset{style=mystyle}

\theoremstyle{definition}
\newtheorem{definition}{Definition}[section]

\theoremstyle{definition}
\newtheorem{exmp}{Example}[section]

\title{Final report: On the Implementability of Global Types}
\author{Student: Gabriele Genovese\\Supervisor: Cinzia Di Giusto}
\date{\today}

\begin{document}

\maketitle

\section*{Abstract}
Behavioral types help define how information is exchanged in distributed systems.  
Multiparty Session Types (MPST) describe interactions between multiple participants  
using global protocols and their local versions. Ensuring correct implementation,  
including deadlock freedom and session conformance, is a key concern in MPST.  
Most research focuses on peer-to-peer communication, but real-world systems use  
different models like mailbox-based or causally ordered messaging. A major challenge  
is that a protocol working in one model may fail in another.  
In this work, we explore a flexible MPST framework that adapts to various communication models,
including unordered, causally ordered, and bus-based messaging. Implementability  
is defined through quasi-synchronous (QS) local types, which ensure deadlock freedom  
and prevent orphan messages while maintaining compatibility with synchronous semantics.  
This approach extends MPST beyond peer-to-peer models, making it useful for  
quasi-synchronous systems, which offer better guarantees for verifying correctness.

\newpage

\tableofcontents

\newpage

\input{content/introduction}

\section{Preliminaries}
In this section, we define the basic knowledge to understand the main contributions
of this work.

\subsection{Multiparty Session Types}

Multiparty Session Types (MPST) provide a formal framework to specify and verify
communication protocols among multiple participants. They ensure that communication
follows a predefined structure, preventing errors such as deadlocks and message loss.
Below, we present the fundamental definitions of MPST. The \textbf{global specification}  
describes the overall communication protocol. From this, we derive the \textbf{local behaviors}  
of each participant through a \textit{projection} operation. The system's \textbf{processes}  
form the \textit{implementation}, defining how participants interact.  
With the definition of a \textit{typing system} and the application of the \textit{type-checking rules},
we ensure that the implementation conforms to the global specification, guaranteeing \textit{well-formedness}.

\subsubsection{The Session Types}

\paragraph{Basic Types}
Basic types are defined as:
\begin{equation}
    T ::= nat \mid int \mid bool
\end{equation}

\paragraph{Global Types}
A global type $G$ describes the overall communication protocol, specifying the
sequence of message exchanges among participants. It defines the expected
interactions at a high level and is formally with the following grammar:

\begin{equation}
G ::= p \to q : \{ l_i(T_i) \}_{i \in I} \mid \mu t. G \mid t \mid G;G \mid \text{end}
\end{equation}

where:
\begin{itemize}
    \item $p, q$ are participants in the communication,
    \item $l_i$ are message labels representing different choices of interaction,
    \item $T_i$ are the types of messages being exchanged,
    \item $\mu t. G$ represents recursion, allowing repetitive behavior,
    \item $t$ is a recursion variable,
    \item $G;G$ denotes sequencing of interactions, meaning that the interactions
    occur in a specified order.
    \item \verb|end| denotes the termination of the communication; it can be omitted.
\end{itemize}

A well-formed global type ensures that all participants can complete their interactions
without deadlocks or undefined behavior.

\paragraph{Local Types}
A local type $L$ represents the communication behavior of an individual participant,
derived from the global type. It specifies the actions that a participant can perform
within the session and is defined as:

\begin{equation}
L ::=  \&_{i \in I}p ? l_i(T_i) \mid \oplus_{i \in I}q ! l_i(T_i) \mid \mu t. L \mid t \mid L;L \mid \text{end}
\end{equation}

where:
\begin{itemize}
    \item $?\{l_i(T_i)\}$ denotes receiving a message labeled $l_i$ with type $T_i$,
    \item $!\{l_i(T_i)\}$ denotes sending a message labeled $l_i$ with type $T_i$,
    \item $\mu t. L$ represents recursion, allowing repetition of behavior,
    \item $t$ is a recursion variable,
    \item $L;L$ defines sequencing of interactions, ensuring messages are exchanged
    in a specified order.
    \item \verb|end| denotes the termination of the communication; it can be omitted.
\end{itemize}

Each participant's local type must be consistent with the global type to ensure
correct communication. Recursion is always guarded.

\paragraph{Projection}
The projection of a global type onto a participant $r$ (where $r \in$
\textit{participants}) is defined by recursion on $G$:
\begin{itemize}
    \item $end \upharpoonright r = end$;
    \item $t \upharpoonright r = t$;
    \item $(\mu t. G) \upharpoonright r = (\mu t. G\upharpoonright r)$;
    \item $G_i;G_j \upharpoonright r = G_i \upharpoonright r \sqcap G_j \upharpoonright r$ where $i, j \in I$;
    \item $p \to r : \{ l_i(T_i) \}_{i \in I} \upharpoonright r = \&_{i \in I}p ? l_i(T_i)$;
    \item $r \to q : \{ l_i(T_i) \}_{i \in I} \upharpoonright r = \oplus_{i \in I}q ! l_i(T_i)$;
    \item $p \to q : \{ l_i(T_i) \}_{i \in I} \upharpoonright r = 0$;
    \item undefined otherwise.
\end{itemize}

$\sqcap$ is the merging operator, that is a partial operation over session types defined as

\begin{equation}
    T1 \sqcap T2 = \begin{cases}
        T1 & \text{if } T1 = T2\\
        T3 & \text{if }\exists I, J:
        \begin{cases}
            T1 = \&_{i \in I} p ? l_i(T_i) & \text{and}\\
            T2 = \&_{j \in J} p ? l_j(T_j)  & \text{and}  \\
            T3 = \&_{k\in I\cup J} p ? l_k(T_k) 
        \end{cases}
        \\
        \text{undefined otherwise}
    \end{cases}
\end{equation}


\paragraph{Well-Formedness}

\begin{theorem}[Subject Reduction]
    Let \(\vdash \mathcal{M} : G\). For all \(\mathcal{M}'\), if \(\mathcal{M} \longrightarrow \mathcal{M}'\),
    then \(\vdash \mathcal{M}' : G'\) for some \(G'\) such that \(G \Longrightarrow G'\).
\end{theorem}

\begin{theorem}[Progress]
    If \(\vdash \mathcal{M} : G\), then either \(\mathcal{M} \equiv p \triangleleft \mathbf{0}\) 
    or there is \(\mathcal{M}'\) such that \(\mathcal{M} \longrightarrow \mathcal{M}'\).
\end{theorem}

As a consequence of subject reduction and progress, we get the safety property stating that a typed multiparty session will never get stuck.

\begin{theorem}[Type Safety]
    If \(\vdash \mathcal{M} : G\), then it does not hold \(\text{stuck}(\mathcal{M})\).
\end{theorem}

These theorems establish the foundation for reasoning about safe and structured
concurrent communications using MPST, ensuring that multi-party protocols
are both reliable and verifiable.

\subsubsection{The Calculus' Syntax}
The process' implementation of the synchronous multiparty session calculus are defined with:

\begin{equation}
    P ::= p ! l \langle e \rangle \mid \sum_{i \in I} p ? l_i (x_i)
    \mid \text{if } e \text{ then } P \text{ else } Q \mid P1;P2 \mid \mu X.P \mid X \mid 0
\end{equation}

\begin{itemize}
    \item $p ! l \langle e \rangle$:  it sends a value of an expression $e$ with label $l$;
    \item $\sum_{i \in I} p ? l_i (x_i)$ is the external choice: the process can receive
    a message with label $l_i$ from process $p$, $x_i$ is the variable that binds the new value received;
    \item $\text{if } e \text{ then } P \text{ else } Q$ represents the internal choice;
    \item $P1;P2$: sequencing the process;
    \item $\mu X.P$: it's a recursive process (we assume they are guarded);
    \item $0$: termination, can be omitted.
\end{itemize}

We define a multiparty session as a parallel composition of pairs (denoted by
$p \triangleleft P$) of participants and processes:

\begin{equation}
 M ::= p \triangleleft P \mid\mid M \mid M
\end{equation}

with the idea that process $P$ represents participant $p$ and interacts with other
processes fulfilling different roles in $M$, a multiparty session is considered
well-formed if all participants are distinct.

\paragraph{Operational semantics}
In the Table~\ref{tab:struc} and \ref{tab:redu}, the operational semantic for the syntax is showed.

\begin{table}[h]
    \centering
    \begin{align*}
        \text{[S-REC]} \quad & \mu X.P \equiv P\{\mu X.P/X\} \\
        \text{[S-MULTI]} \quad & P \equiv Q \Rightarrow p \triangleleft P \mid M \equiv p \triangleleft Q \mid M \\
        \text{[S-PAR 1]} \quad & p \triangleleft \mathbf{0} \mid M \equiv M \\
        \text{[S-PAR 2]} \quad & M \mid M' \equiv M' \mid M \\
        \text{[S-PAR 3]} \quad & (M \mid M') \mid M'' \equiv M \mid (M' \mid M'')
    \end{align*}
    \caption{Structural congruence.}
    \label{tab:struc}
\end{table}

\begin{table}[h]
    \centering
    \begin{align*}
        \text{[R-COMM]} \quad & \frac{j \in I \quad e \downarrow v}{p \triangleleft \sum_{i \in I} q?l_i(x).P_i \mid q \triangleleft p!l_j\langle e \rangle.Q \mid M \longrightarrow p \triangleleft P_j\{v/x\} \mid q \triangleleft Q \mid M} \\[10pt]
        \text{[T-CONDITIONAL]} \quad & \frac{e \downarrow \text{true}}{p \triangleleft \text{if } e \text{ then } P \text{ else } Q \mid M \longrightarrow p \triangleleft P \mid M} \\[10pt]
        \text{[F-CONDITIONAL]} \quad & \frac{e \downarrow \text{false}}{p \triangleleft \text{if } e \text{ then } P \text{ else } Q \mid M \longrightarrow p \triangleleft Q \mid M} \\[10pt]
        \text{[R-STRUCT]} \quad & \frac{M_1' \equiv M_1 \quad M_1 \longrightarrow M_2 \quad M_2 \equiv M_2'}{M_1' \longrightarrow M_2'}
    \end{align*}
    \caption{Reduction rules.}
    \label{tab:redu}
\end{table}

\paragraph{Type system}
We now introduce a type system for the multiparty session calculus presented. 
We distinguish three kinds of typing judgments:
\begin{itemize}
    \item $\Gamma \vdash e : S$
    \item $\Gamma \vdash P : S$
    \item $ \vdash M : G$
\end{itemize}
where $\Gamma$ (called the context) is the \textit{typing environment}:
\begin{equation}
    \Gamma ::= \empty \mid\mid \Gamma, x : S \mid\mid \Gamma, X : T
\end{equation}
i.e., a mapping that associates expression variables with sorts, and process variables with session types.
We say that a multiparty session M is well typed if there is a global type G such that ` M : G.
We omit the rules for the expression because they are trivial. In Table~\ref{tab:typ}, there are
the non-trivial rules.

\begin{table}[h]
    \centering
    \begin{align*}
        \text{[T-SUB]} \quad & \frac{\Gamma \vdash P : T \quad T \leq T'}{\Gamma \vdash P : T'} \\[10pt]
        \text{[T-0]} \quad & \Gamma \vdash \mathbf{0} : \text{end} \\[10pt]
        \text{[T-REC]} \quad & \frac{\Gamma, X : T \vdash P : T}{\Gamma \vdash \mu X.P : T} \\[10pt]
        \text{[T-VAR]} \quad & \Gamma, X : T \vdash X : T \\[10pt]
        \text{[T-INPUT-CHOICE]} \quad & 
        \frac{\forall i \in I \quad \Gamma, x_i : S_i \vdash P_i : T_i}
        {\Gamma \vdash \sum_{i \in I} q?l_i(x_i).P_i : \&_{i \in I} q?l_i(S_i).T_i} \\[10pt]
        \text{[T-OUT]} \quad & 
        \frac{\Gamma \vdash e : S \quad \Gamma \vdash P : T}
        {\Gamma \vdash q!l\langle e \rangle.P : q!l(S).T} \\[10pt]
        \text{[T-CHOICE]} \quad & 
        \frac{\Gamma \vdash e : \text{bool} \quad \Gamma \vdash P_1 : T \quad \Gamma \vdash P_2 : T}
        {\Gamma \vdash \text{if } e \text{ then } P_1 \text{ else } P_2 : T} \\[10pt]
        \text{[T-SESS]} \quad & 
        \frac{\forall i \in I \quad \vdash P_i : G[p_i] \quad pt\{G\} \subseteq \{p_i \mid i \in I\}}
        {\vdash \prod_{i \in I} p_i \triangleleft P_i : G}
    \end{align*}
    \caption{Typing rules for processes and sessions.}
    \label{tab:typ}
\end{table}

\subsubsection{Example}


\paragraph{Global Specification}
The global communication protocol is defined as:
\begin{align*}
    &\text{Customer} \to \text{Agency} : \text{details}(\mathbb{N}). \\
    &\text{Agency} \to \text{Hotel} : \text{details}(\mathbb{N}). \\
    &\text{Hotel} \to \text{Customer} : \text{ok}(\mathbb{B}). \text{end}
\end{align*}

\paragraph{Local Types}
\begin{align*}
    T_C &= \text{Agency}! \text{details}(\mathbb{N}). \text{Hotel}? \text{ok}(\mathbb{B}). \text{end} \\
    T_A &= \text{Customer}? \text{details}(\mathbb{N}). \text{Hotel}! \text{details}(\mathbb{N}). \text{end} \\
    T_H &= \text{Agency}? \text{details}(\mathbb{N}). \text{Customer}! \text{ok}(\mathbb{B}). \text{end}
\end{align*}

\paragraph{Process Implementation}
\begin{align*}
    P_C &= \text{Agency}! \text{details}(42). \text{Hotel}? \text{ok}(b). \mathbf{0} \\
    P_A &= \text{Customer}? \text{details}(d). \text{Hotel}! \text{details}(d). \mathbf{0} \\
    P_H &= \text{Agency}? \text{details}(d). \text{Customer}! \text{ok}(\text{true}). \mathbf{0}
\end{align*}

\paragraph{Typing Proof}
We prove that each process conforms to its local type.

\begin{lemma}[Customer Typing]
If \(\Gamma \vdash P_C : T_C\), then \(P_C\) follows \(T_C\).
\begin{proof}
By applying the \textbf{T-OUT} rule:
\[
\frac{\Gamma \vdash 42 : \mathbb{N} \quad \Gamma \vdash P_C' : \text{Hotel}? \text{ok}(\mathbb{B}). \text{end}}
{\Gamma \vdash \text{Agency}! \text{details}(42). P_C' : \text{Agency}! \text{details}(\mathbb{N}). T_C'}
\]
where \(P_C' = \text{Hotel}? \text{ok}(b). \mathbf{0}\).

Applying \textbf{T-INPUT-CHOICE}:
\[
\frac{\Gamma \vdash b : \mathbb{B} \quad \Gamma \vdash \mathbf{0} : \text{end}}
{\Gamma \vdash P_C' : \text{Hotel}? \text{ok}(\mathbb{B}). \text{end}}
\]
Thus, \(P_C\) is well-typed.
\end{proof}
\end{lemma}

\begin{lemma}[Agency Typing]
\begin{proof}
Applying \textbf{T-INPUT-CHOICE} and \textbf{T-OUT}:
\[
\frac{\Gamma \vdash d : \mathbb{N} \quad \Gamma \vdash \mathbf{0} : \text{end}}
{\Gamma \vdash P_A : \text{Customer}? \text{details}(\mathbb{N}). \text{Hotel}! \text{details}(\mathbb{N}). \text{end}}
\]
\end{proof}
\end{lemma}

\begin{lemma}[Hotel Typing]
\begin{proof}
Applying \textbf{T-INPUT-CHOICE} and \textbf{T-OUT}:
\[
\frac{\Gamma \vdash d : \mathbb{N} \quad \Gamma \vdash \mathbf{0} : \text{end}}
{\Gamma \vdash P_H : \text{Agency}? \text{details}(\mathbb{N}). \text{Customer}! \text{ok}(\mathbb{B}). \text{end}}
\]
\end{proof}
\end{lemma}

\begin{theorem}[Session Typing]
\begin{proof}
By \textbf{T-SESS}:
\[
\frac{\forall i \in \{C, A, H\}, \Gamma \vdash P_i : T_i}
{\Gamma \vdash P_C \mid P_A \mid P_H : G}
\]
Since \(T_C, T_A, T_H\) are well-typed, the system follows \(G\).
\end{proof}
\end{theorem}

\paragraph{Operational Semantics Execution}
The execution follows the operational semantics rules:

\begin{align*}
    P_C \mid P_A \mid P_H &\to 
    (\text{Agency}! \text{details}(42). P_C') \mid (\text{Customer}? \text{details}(d). P_A') \mid P_H \\
    &\to P_C' \mid (\text{Hotel}! \text{details}(42). P_A') \mid P_H \quad \text{(Comm. between Customer, Agency)} \\
    &\to P_C' \mid P_A' \mid (\text{Customer}! \text{ok}(\text{true}). P_H') \quad \text{(Comm. between Agency, Hotel)} \\
    &\to P_C' \mid P_A' \mid P_H' \quad \text{(Comm. between Hotel, Customer)}
\end{align*}

At the end of the execution, all processes reach \(\mathbf{0}\), meaning communication terminates successfully.
Thus, the protocol is correctly implemented and respects the MPST specifications.



\section{State of the art}

In this section realizability and implementability will be used as the same problem.
Informally, an MSC can be defined realizable if a distributed 
implementation exists that exactly reproduces the graph's behaviors.

\subsection{Realizability by Alur}

\subsubsection{Definitions}

\begin{definition}[Weakly-implies]
A set $L$ of MSCs weakly implies an MSC $M$ iff for all $1\leq i \leq n$, 
there exists an MSC $M_i \in L$ such that $M|_i = M_i|_i$.
\end{definition}

\begin{definition}{Weakly realizable}
The weak closure $L^w$ of a set $L$ of MSCs contains all the MSCs $L$ weakly implies.
The set $L$ is weakly realizable iff $L=L^w$. The MSC-graph $G$ is said to be weakly
realizable if the set $L(G)$ of MSCs is. Thus, a weakly realizable graph already contains
all the implied scenarios.
\end{definition}

\begin{definition}[Deadlock-free]
Consider a set $A_i$ of concurrent automata and the product $A =\prod_i A_i$
A state $q$ of the product $A$ is said to be a deadlock state if no accepting state of $A$ is reachable from $q$.
The product $A$ is said to be deadlock-free if no state reachable from its initial state is a deadlock state.
\end{definition}

\begin{definition}[Safely realizable]
A set $L$ of MSCs is said to be safely realizable if $L=L\prod A_i$ for some $\langle A_i | 1\leq i \leq n \rangle$
such that $\prod A_i$ is deadlock-free.
\end{definition}

\subsubsection{Realizability problem}

In \cite{alur2005realizability}, they prove many interesting results within the scope of MSC graphs.
They first examine the realizability of bounded MSC-graphs. Two types are considered: 
\textit{weak} (allowing deadlocks) and \textit{safe} (deadlock-free). While weak realizability is \verb|coNP|-complete
and safe realizability solvable in polynomial time for finite MSC sets, they show that for
bounded MSC-graphs, weak realizability becomes undecidable, and safe realizability lies in 
\verb|EXPSPACE|.

\begin{theorem}
Given a bounded MSC graph G, checking if G is weakly realizable is undecidable.
\end{theorem}

\begin{proof}
The proof is a reduction from the post correspondence problem (PCP). In the paper \cite{alur2005realizability}
they further reduce the problem to a relaxed version. First, they reduce PCP to OneSidedPCP and then to
Relaxed PCP (RPCP). We skip this part as is not meaningful for the proof. We can state the RPCP to be the following
problem: given $\{(v_1, w_1), (v_2, w_2), ..., (v_r, w_r)\}$, determinate whether there are indices $i_1, ..., i_m$
such that $x_{i_1}...x_{i_n}=y_{i_1}...y_{i_m}$, where $x_{i_j}, y_{i_j} \in \{v_{i_j},w_{i_j}\}$, for some
index $i_l \  x_{i_l}\neq y_{i_l}$, and for all $j\leq m,y_{i_1}...y_{i_j}$ is a prefix of $x_{i_1}...x_{i_j}$.
The relaxed problem essential says that must have the left side of the equivalence that grows more than the right side,
we must have also some tuple that are different, and that we can choose whatever side of tuple for the composition
of the string. Now we reduce RPCP to weak realizability. 

\newcommand{\syncmess}[3]{%
  \mess{#1}{#2}{#3}%
  \mess{}{#3}{#2}%
  \nextlevel
}

\newcommand{\asyncmess}[3]{%
  \mess{#1}{#2}{#3}%
  \nextlevel
}

\begin{figure}[h]
\centering
\begin{msc}[draw frame=none, draw head=none, msc keyword=, head height=0px, label distance=0.5ex, foot height=0px, foot distance=0px]{}
    \declinst{P1}{P1}{}
    \declinst{P2}{P2}{}
    \declinst{P3}{P3}{}
    \declinst{P4}{P4}{}

    \syncmess{$(i,0)$}{P1}{P2}
    \asyncmess{$i$}{P1}{P4}
    \syncmess{$(i,0)$}{P3}{P4}
    \syncmess{$v_i^1$}{P2}{P3}
    \syncmess{...}{P2}{P3}
    \syncmess{$v_i^c$}{P2}{P3}
\end{msc}
\caption{The $M_i^0$ MSC. For a string $u$, let $u_l$ denote the $l$’th character of the string.}
\end{figure}

\begin{figure}[h]
\centering
\begin{msc}[draw frame=none, draw head=none, msc keyword=, head height=0px, label distance=0.5ex, foot height=0px, foot distance=0px]{}
    \declinst{P1}{P1}{}
    \declinst{P2}{P2}{}
    \declinst{P3}{P3}{}
    \declinst{P4}{P4}{}

    \syncmess{$(i,1)$}{P1}{P2}
    \asyncmess{$i$}{P1}{P4}
    \syncmess{$(i,1)$}{P3}{P4}
    \syncmess{$w_i^1$}{P2}{P3}
    \syncmess{...}{P2}{P3}
    \syncmess{$w_i^d$}{P2}{P3}
\end{msc}
\caption{The $M_i^1$ MSC. For a string $u$, let $u_l$ denote the $l$’th character of the string.}
\end{figure}


\textit{Encodig:} Given a finite set $L$ of
MSCs, let $L^*$ denote the MSC graph that consists of the complete
graph with $|L|$ vertices, one per MSC in the set $L$, dummy initial
and terminal vertices $v_I$, $v_T$ with empty MSCs, and edges from
$v_I$ to all vertices of $L$ and from those to $v_T$. Thus, an MSC of
this graph is simply a concatenation of MSCs from the set $L$.

In the sequel, we say that a process $p$ synchronously sends a
message $m$ to process $q$, if $p$ sends $m$ to $q$ immediately
followed by $q$ sending the message $m$ back to $p$. In figures,
such messages will be depicted by double arrows.

Given an instance $\Delta = \{(v_1, w_1), \ldots, (v_m, w_m)\}$ of
RPCP, we build a set $L$ of MSCs over 4 processes as follows.
For each pair $(v_i, w_i)$ we build two MSCs $M^0_i$ and $M^1_i$,
which are depicted in Fig.~3.
Thus in $M^0_i$, process 1 sends synchronously $(i, 0)$ to process 2,
then sends the index $i$ to process 4, and then process 4 sends
synchronously $(i, 0)$ to process 3. After that, process 2
synchronously sends the sequence of characters of $v_i$ to process 3 (note we assume $c$
is the length of $v_i$ and $d$ the length of $w_i$ in the figure), $M^1_i$
is similar. Observe that the communication graph of each of these
MSCs is strongly connected and involves all the processes, and
hence, the MSC graph $L^*$ is bounded.

\textbf{Claim 1:} $\Delta \in \text{RPCP}$ iff $L^*$ is not weakly realizable.

\textit{Proof.} For the ``only if" direction, suppose $R = (i_1, a_1, b_1, i_2, a_2, b_2, \ldots,
i_m, a_m, b_m)$ are the indices for a solution to $\mathcal{B}$, and the bits $a_j$ and $b_j$
indicate which string ($v_{i_j}$ or $w_{i_j}$) is chosen to go into the two (left and right)
long strings.

Consider the new MSCs $M$ and $M'$ obtained from the sequences
$M = M^{a_1}_{i_1} \cdots M^{a_m}_{i_m}$ and $M' = M^{b_1}_{i_1} \cdots M^{b_m}_{i_m}$.
Executions of both of these (sequences of) MSCs must exist in any
realization of $L^*$. We then look at the projections $M|_1$, $M|_2$, $M|_3$,
and $M|_4$ of $M$, and $M'|_1$, $M'|_2$, $M'|_3$, $M'|_4$ of $M'$ onto the
4 processes. Now consider an MSC $M''$ formed from $M'|_1$, $M'|_2$,
$M|_3$, and $M|_4$. The claim is that the combined MSC $M''$ is weakly
implied by $L^*$. By definition, the only thing to establish is that $M''$
is indeed an MSC, in the sense that it is acyclic, well-formed, and complete.

The only new situation in terms of communication in $M''$ is the
communication between $P_1$ and $P_4$, and between $P_2$ and $P_3$.
But the communication between $P_1$ and $P_4$ is consistent in
$M'|_1$ and $M|_4$ (i.e., the sequence of messages sent from $P_1$ to
$P_4$ in $M'|_1$ is equal to the sequence of messages received in $M|_4$),
and the communication between $P_2$ and $P_3$ is consistent in
$M'|_2$ and $M|_3$ because $R$ is a solution to the RPCP.

Furthermore, the acyclicity of $M''$ follows from the property of the
solution that the string formed by the first $j$ words on processes 1
and 2 is always a prefix of the string formed by the first $j$ words
on processes 3 and 4. Consequently, each message from $P_1$ to $P_4$
is sent before it needs to be received.

But note that $M''$ cannot itself be in $L^*$ because there must be
some index $i_j$ where $a_j \neq b_j$, and no MSC exists in $L$ where,
after process 1 announces the index, what process 2 sends is not
identical to what process 3 receives.

Now, for the “if” direction, suppose there is some MSC $M^@$ which
exists in any realization of $L^*$, but is not in $L^*$ itself. We want
to derive a solution to $\mathcal{B}$ from $M^@$.

First, it is clear that the projection $M^@|_1$ must consist of a sequence
of pairs of messages (the first of each pair acknowledged), sent from
process 1 to processes 2 and 4, respectively, with messages $(i, b)$ and $i$.

Likewise, it is clear that, in order for process 2 to receive those messages,
$M^@|_2$ must consist of a sequence of receipts of $(i, b)$ pairs, and after
each $(i, b)$, either $v_i$ or $w_i$ is sent to process 3, based on whether
$b = 0$ or $b = 1$, before the next index pair is received.

Likewise, $M^@|_4$ consists of a sequence of receipts of index $i$ from
process 1, followed by sending of $(i, 0)$ or $(i, 1)$ to process 3, and
$M^@|_3$ consists of a sequence of receipt of $(i, 0)$ or $(i, 1)$ followed
by receipt of $v_i$ or $w_i$, respectively.

Now, since $M^@$ is not in $L^*$, for some index $i$ the choice of $v_i$ or
$w_i$ must differ on process 2 and process 3. (Note, we are assuming that
the buffers between processes are FIFO.)

Furthermore, because of the precedences, the prefix formed by the first
$j$ words on process 2 must precede the $(j + 1)$-th message from
process 1 to process 4, which in turn precedes the $(j + 1)$-th message
from 4 to 3, and hence the $(j + 1)$-th word on process 3. That is, the
string formed by the first $j$ words on process 2 is a prefix of the string
formed by the first $j$ words on process 3. Therefore, we can readily
build a solution for $\mathcal{B}$ from $M^@$.

\end{proof}

\subsection{Implementability by Stutz}

Part taken from \cite{stutz2024implementability}.
To start with, let $P$ be a finite set of participants and $V$ be a finite set of messages
and $\Gamma$ the set of all send and receive events. $\Gamma_p$ is the set of all send and receive events
regarding a participant $p\in P$.

\subsubsection{Definitions}

\newcommand{\send}[3]{%
  #1\rhd #2 ! #3
}

\newcommand{\recv}[3]{%
  #1\lhd #2 ? #3
}

\newcommand{\Epsilon}[0]{%
  \mathcal{E}
}

\newcommand{\doublebrace}[1]{%
  \{\!\!\{#1\}\!\!\}
}

\begin{definition}{CSM}
A \textit{communicating state machine} (CSM) is $\mathsf{A} = \doublebrace{A_p}_{p\in P}$
over $P$ and $V$ if $A_p$ is a finite state machine over $\Gamma_p$ for every $p \in P$,
denoted by $(Q_p,\Gamma_p,\delta_p,q_{0,p}, F_p)$. Let $\prod_{p\in P} Q_p$ denote
the set of global states and \verb|Chan| $=\{(p,q)| p,q \in P,p\neq q\}$ denote the set
of channels. A configuration of $A$ is a pair $(\overrightarrow{q}, \Epsilon)$, where $\overrightarrow{q}$ is a vector
of sates, one for each participant, and $\Epsilon:\text{Chan}\to V^*$ is a mapping from 
each channel to a sequence of messages. We may write $\Epsilon(p, q)$ for $\Epsilon((p, q))$.
We use $\overrightarrow{q}_p$ to denote the state of $p$ in $\overrightarrow{q}$. The
relation, denoted $\to$, is defined as follows:
\begin{itemize}
  \item $(\overrightarrow{q}, \Epsilon) \xrightarrow{\send{p}{q}{m}} (\overrightarrow{q}', \Epsilon')$
  if $(\overrightarrow{q_p},\send{p}{q}{m},\overrightarrow{q_p}') \in \delta_p, \overrightarrow{q_r} =\overrightarrow{q_r}'$
  for every participant $r\neq p$, $\Epsilon'((p,q)) = \Epsilon((p,q))\cdot m$ and $\Epsilon'(c) = \Epsilon(c)$
  for every other channel $c \in $ Chan.
  \item $(\overrightarrow{q}, \Epsilon) \xrightarrow{\recv{q}{p}{m}} (\overrightarrow{q}', \Epsilon')$
  if $(\overrightarrow{q_p},\recv{q}{p}{m},\overrightarrow{q_p}') \in \delta_p, \overrightarrow{q_r} =\overrightarrow{q_r}'$
  for every participant $r\neq p$, $\Epsilon((p,q)) = m\cdot\Epsilon'((p,q))$ and $\Epsilon'(c) = \Epsilon(c)$
  for every other channel $c \in $ Chan.
  \item $(\overrightarrow{q}, \Epsilon) \xrightarrow{\epsilon} (\overrightarrow{q}', \Epsilon)$ if 
  $(\overrightarrow{q_p},\epsilon,\overrightarrow{q_p}') \in \delta_p$ for some participant $p$, and
  $\overrightarrow{q_q} =\overrightarrow{q_q}'$ for every participant $q\neq p$.
\end{itemize}
\end{definition}

Informally, a configuration $(\overrightarrow{q}, \Epsilon)$ is a deadlock if it is not final and has no outgoing 
transitions while it is reachable if it occurs on some run of A. A CSM is deadlock-free if no reachable
configuration is a deadlock. $\Gamma^\omega$ denotes the infinite set of events.

\begin{definition}[Implementability problem]
A language $L \subseteq \Gamma^\omega$ is said to be implementable if there exists a CSM $\doublebrace{A_p}_{p\in P}$ such that
\begin{itemize}
  \item deadlock freedom: $\doublebrace{A_p}_{p\in P}$ is deadlock-free, and
  \item protocol fidelity: $L$ is the language of $\doublebrace{A_p}_{p\in P}$.
\end{itemize}
\end{definition}



\section{Content}

\subsection{Useful definitions}

In our setting, global types are automata that describe a language of MSCs. 
An \emph{arrow} is a triple $(p,q,m)\in\Procs\times\Procs\times\Messages$ with $p\ne q$;
we often write $\marrow{p}{q}{m}$ instead of $(p,q,m)$, and write $\labelalphabet$
to denote the finite set of arrows.

\begin{definition}[Global Type]
A global type $\gt$ is a deterministic finite state automaton over the alphabet $\labelalphabet$.
\end{definition}

Informally, weakly-realizable means that it's realizable and deadlock is allowed.
%% definire da linguaggio di G
\begin{definition}[Weakly-realizable]
It's the 
weak closure $L^w$ of a set $L$ of global type contains all the global type $L$ weakly implies.
The set $L$ is weakly realizable iff $L=L^w$. The global type-graph $G$ is said to be weakly
realizable if the set $L(G)$ of global type is. Thus, a weakly realizable graph already contains
all the implied scenarios.
\end{definition}


\begin{definition}[Relaxed Post Correspondence Problem (RPCP) problem]
Given $\{(v_1, w_1), (v_2, w_2), ..., (v_r, w_r)\}$, determinate whether there are indices $i_1, ..., i_m$
such that $x_{i_1}...x_{i_n}=y_{i_1}...y_{i_m}$, where $x_{i_j}, y_{i_j} \in \{v_{i_j},w_{i_j}\}$, for some
index $i_l \  x_{i_l}\neq y_{i_l}$, and for all $j\leq m,y_{i_1}...y_{i_j}$ is a prefix of $x_{i_1}...x_{i_j}$.
The relaxed problem essentially says that must have the left side of the equivalence that grows more than the right side,
we must have also some tuple that are different, and that we can choose whatever side of tuple for the composition
of the string.
\end{definition}

RPCP is proven to be an undecidable problem by reduction from the classic
PCP problem by Alur, et al. \cite{alur2005realizability}.

\subsection{Proof}

\begin{theorem}
Given a global type $G$, checking if $G$ is weakly-realizable is undecidable.
\end{theorem}

\begin{proof}
The proof proceeds via a reduction from the RPCP problem. Let's define some useful
elements for the proof.

Given a finite set $L$ of global types, where each is defined as a deterministic
finite state automaton over the alphabet $\labelalphabet$, we define $L^*$
to be the language consisting of all finite concatenations of languages in $L$.
To construct an automaton recognizing $L^*$, we introduce an initial state
$v_I$ and a terminal state $v_T$, and for each automaton in $L$, we add
a transition from $v_I$ to its initial state and from each of its final states
to $v_T$. Additionally, we connect every final state of each automaton in $L$
to the initial state of every other automaton in $L$. Thus, $L^*$ contains all
MSCs that can be obtained by concatenating MSCs from individual languages in $L$.

Given an instance $\Delta = {(v_1, w_1), \ldots, (v_m, w_m)}$ of RPCP,
we build a set $L$ of languages over 4 processes as follows.
For each pair $(v_i, w_i)$ we construct a language $\mathcal{L}_i^n$
(where $n \in {0,1}$) consisting of MSCs specified by an automaton
$A_i^n$ (described in Fig.~\ref{fig:gtype}).

In each global type $G_i^n$, process $q$ sends a message $(i,n)$ to itself,
then sends the index $i$ to process 4; process 4 then sends $(i,n)$
to process 3. After this, process 2 sends the characters of $v_i$
(or $w_i$, depending on $n$) to process 3. This communication pattern
ensures that the MSCs in $\mathcal{L}_i^n$ describe structured
interactions with explicit acknowledgments.
Observe that the communication graph of each MSC in $\mathcal{L}_i^n$
is strongly connected and involves all processes.

\begin{figure}[h]
\centering
\begin{tikzpicture}[->, node distance=24mm, on grid, auto]
  \node[state] (q0) {$q_0$};
  \node[state] (q1) [right=of q0] {$q_1$};
  \node[state] (q2) [right=of q1] {$q_2$};
  \node[state] (q3) [right=of q2] {$q_3$};
  \node[state] (q4) [right=of q3] {$q_4$};
  \node[state] (q5) [right=of q4] {$q_x$};
  \node[state] (q6) [right=of q5] {$q_y$};

  \path (q0) edge[] node[above] {\arrmess{p}{q}{(i,n)}} (q1);
  \path (q1) edge[] node[above] {\arrmess{p}{s}{i}} (q2);
  \path (q2) edge[] node[above] {\arrmess{s}{r}{(i,n)}} (q3);
  \path (q3) edge[] node[above] {\arrmess{q}{r}{x_i^1}} (q4);
  \path (q4) edge[] node[above] {\arrmess{q}{r}{...}} (q5);
  \path (q5) edge[] node[above] {\arrmess{q}{r}{x_i^c}} (q6);
\end{tikzpicture}
\caption{The global type $G^n_i$ where $n\in\{0,1\}$ and $x\in\{v,w\}$.}
\label{fig:gtype}
\end{figure}

We need to prove:
\begin{center}
$\Delta \in \text{RPCP}$ if and only if the language $L^*$ is not realizable.
\end{center}

\begin{itemize}

  \item[$\Rightarrow$] Suppose there is a solution to RPCP: $\Delta = \{i_1,\ldots,i_k$\}
  with $v_{i_1}\ldots v_{i_k} = w_{i_1}\ldots w_{i_k}$. Then, ...
  
  \item[$\Leftarrow$] If the language $L^*$ is not realizable,...

\end{itemize}

\textbf{Lemma.} Each $\mathcal{L}_i^n$ is a synch-language.

\begin{proof}
We show that the MSCs in $\mathcal{L}_i^n$ admit only synchronous linearizations.
Due to the presence of ACK messages, send and receive
events are closely paired. The only possible reordering involves
the ACK from \arrmess{p}{s}{i}, ...
\end{proof}


\end{proof}




\newpage

\bibliographystyle{plain}
\bibliography{ref}

\end{document}